\begin{dashboardbox}[投资结论]
\noindent\begin{tabularx}{\textwidth}{L{3.0cm}X L{3.0cm}X}
\textbf{股票} & 东阳光（600673.SH） & \textbf{当前价} & 38.99元 \\
\textbf{评级} & 中性／持有／事件驱动观察 & \textbf{最终目标} & 38.35元 \\
\textbf{基本面锚} & 31.26元 & \textbf{隐含空间} & -1.6\% \\
\textbf{情景区间} & 22.10--48.87元 & \textbf{主要风险} & 稀释、融资、验收、客户集中、商誉 \\
\end{tabularx}
\end{dashboardbox}

\begin{houseviewbox}[本机构观点]
7月10日新增130--150亿元C单，且公司首次明确前序合同已交付、验收并确认收入，经营证据明显增强；股东会也已通过秦淮交易。但旧43元目标忽略40905万股购买资产发股、潜在配套融资稀释，并将秦淮期权估得远高于官方交易权益价值。重算后，基本面锚为31.26元，市场调整目标为38.35元。结论不是追高买入，而是已有仓位持有、等待监管与验收硬证据。
\end{houseviewbox}

\begin{exhibitbox}[三情景估值]
\noindent\begin{tabularx}{\textwidth}{L{1.6cm}R{1.5cm}R{1.4cm}R{1.4cm}R{1.2cm}R{1.2cm}X}
\toprule
\textbf{情景} & \textbf{净利润} & \textbf{股本} & \textbf{EPS} & \textbf{PE} & \textbf{价值} & \textbf{核心假设} \\
\midrule
bear & 19.00 & 30.10 & 0.631 & 35x & 22.10 & 交易完成延后，公司仅保留现有30\%经济权益；算力合同不计可靠全年利润，传统主业低于券商基准。 \\
base & 26.91 & 36.75 & 0.732 & 45x & 32.95 & 主业达到国金最新原始PDF预测；秦淮采用2025备考归母增量；三单年化未税收入的25\%在2026年确认，净利率4\%；计入购买资产发股和按监管价格下限完成的80亿元配套融资。 \\
bull & 32.40 & 36.47 & 0.888 & 55x & 48.87 & 材料主业超预期、秦淮整合顺利；三单年化未税收入40\%在2026年确认，净利率7.5\%；配套融资价格为现价90\%。 \\
\bottomrule
\end{tabularx}
\end{exhibitbox}

\section{行动框架}

新开仓不建议在39元附近追涨。已有仓位若未触发基本面失效可持有；37.15--37.32元是事件缺口观察位，34.77--35元是深度支撑，45.16元以上必须由监管推进、配套融资定价或C单验收支撑。

\sourcenote{实时行情包；三份算力合同公告；重组草案；国金原始研报；data/current\_valuation\_model\_20260713.json。}
